\documentclass[aps,prb,preprint,superscriptaddress,nofootinbib]{revtex4-2}
\usepackage[T1]{fontenc}
\usepackage{amsmath,mathtools}
\usepackage{newtxtext,newtxmath}
\usepackage{bm}
\usepackage{booktabs}
\usepackage{microtype}
\usepackage{xcolor}
\usepackage{hyperref}
\usepackage[nameinlink,capitalise]{cleveref}
\hypersetup{
  colorlinks=true,
  linkcolor=blue!45!black,
  citecolor=blue!45!black,
  urlcolor=blue!45!black,
  pdftitle={Wannier-Green-Function Theory of Displacement-Induced Spin Torque for Magnon-Polaron Coupling},
  pdfauthor={Author Name},
  pdfsubject={Relativistic spin-lattice response in a spinor Wannier basis},
  pdfkeywords={LKAG, spin torque, DFPT, Wannier functions, magnon-polaron}
}

\newcommand{\kk}{\mathbf k}
\newcommand{\qq}{\mathbf q}
\newcommand{\rr}{\mathbf R}
\newcommand{\Tr}{\operatorname{Tr}}
\newcommand{\tr}{\operatorname{tr}}
\newcommand{\ii}{\mathrm{i}}
\newcommand{\dd}{\mathrm{d}}
\newcommand{\XC}{\mathrm{XC}}
\newcommand{\TRS}{\mathrm{TRS}}
\newcommand{\bA}{\mathcal A}
\newcommand{\bF}{\mathcal F}
\newcommand{\bT}{\mathcal T}
\newcommand{\bL}{\mathcal L}
\newcommand{\AuthorName}{Author Name}

\begin{document}

\title{Wannier-Green-Function Theory of Displacement-Induced Spin Torque for Magnon-Polaron Coupling}
\author{\AuthorName}
\affiliation{Affiliation to be inserted}
\date{\today}

\begin{abstract}
We formulate a q-resolved mixed spin-lattice response in an orthonormal spinor Wannier basis for direct use in magnon-polaron calculations. The construction combines two established ingredients: a relativistic localized-orbital LKAG torque formalism in which only the exchange field is rotated, and a spin-lattice magnetic-force-theorem expansion whose site-diagonal term has the structure of a spin vertex, two propagators, and a lattice perturbation. The lattice vertex is supplied by a screened density-functional perturbation theory Hamiltonian derivative, while spin-orbit coupling remains in the full electronic Green function. We define a Hermitian local partition of the exchange field, distinguish transverse-direction from rotation-angle coordinates, derive a finite-q bubble and its retarded/advanced q-pair completion, and state the optional direct mixed vertex required for a total derivative. The resulting Cartesian torque kernel is projected onto phonon zero-point displacements and externally computed paraunitary magnon eigenvectors. Source-derived statements and project extensions are identified explicitly, and a validation hierarchy is given for time reversal, locality, finite differences, gauge covariance, acoustic translations, and bosonic rephasing.
\end{abstract}

\maketitle
\tableofcontents

\section{Introduction}
\label{sec:intro}

The Liechtenstein-Katsnelson-Antropov-Gubanov (LKAG) construction maps infinitesimal electronic spin rotations onto parameters of an effective spin Hamiltonian using the magnetic force theorem. Its nonrelativistic and relativistic forms are now standard tools for exchange, anisotropic exchange, and Dzyaloshinskii-Moriya interactions~\cite{Liechtenstein1987,Udvardi2003,EbertMankovsky2009,Szilva2023}. Localized-orbital implementations are attractive because the magnetic degrees of freedom are already represented in a compact basis, but the definition of a local operator and the choice of the rotated magnetic subspace require care~\cite{Oroszlany2019,SorianoPalacios2014,He2021TB2J}.

Martinez-Carracedo \emph{et al.} developed a fully relativistic LKAG mapping for nonorthogonal pseudoatomic orbitals, proposed a Hermitian local projection, and demonstrated that a faithful DFT-to-spin mapping rotates the exchange field rather than the complete spinor Hamiltonian~\cite{MartinezCarracedo2023}. Independently, Mankovsky \emph{et al.} treated spin tilts and atomic displacements as perturbations on equal footing in relativistic multiple scattering. Their site-diagonal spin-lattice coefficient contains a torque vertex, a displacement vertex, and two scattering propagators, and it produces a displacement-induced effective magnetic field and torque~\cite{Mankovsky2023SLC}.

The goal here is not to rederive either framework in its native basis. Instead, we formulate their common perturbative structure in an orthonormal spinor Wannier representation and couple it to a screened DFPT perturbation. Wannier interpolation provides a natural basis for both the Hamiltonian and electron-phonon matrix elements~\cite{Marzari2012,Giustino2007,Pizzi2020}, while DFPT supplies the screened first-order lattice response~\cite{Baroni2001}. The electronic calculation outputs a transverse-spin/displacement Hessian. Phonon and magnon spectra are assumed to have been computed independently, and the electronic kernel is projected onto their eigenvectors.

The intended application is a linear one-magnon--one-phonon hybridization in a spin-orbit-coupled magnet. Atomistic spin-lattice Hamiltonians and coupled spin-lattice dynamics provide the broader modeling context. The present linear torque channel differs from a scalar exchange-striction derivative around a collinear state, which commonly produces two-magnon--one-phonon interactions. A total electronic torque response can contain local anisotropy, DMI-like, and symmetric anisotropic intersite contributions without resolving them separately. The direct mode projection therefore avoids an unnecessary spin-model decomposition before constructing the magnon-polaron Hamiltonian.

\section{Scope, assumptions, and provenance}
\label{sec:scope}

We distinguish three levels of provenance throughout.

\begin{enumerate}
\item \emph{Source-derived ingredients.} The local projection, exchange-field-only rotation, and commutator torque vertex follow Ref.~\cite{MartinezCarracedo2023}. The mixed spin-displacement Green-function structure, effective-field interpretation, and phonon-like finite-q motivation follow Ref.~\cite{Mankovsky2023SLC}.
\item \emph{Standard supporting formalism.} DFPT, Wannier interpolation, and bosonic paraunitary diagonalization follow Refs.~\cite{Baroni2001,Giustino2007,Marzari2012,Pizzi2020,Colpa1978}.
\item \emph{Project extensions.} The screened-DFPT Wannier substitution, arbitrary-gauge time-reversal sewing contract, explicit finite-q q-pair completion, and projection onto external magnon modes are the formulation developed for this implementation.
\end{enumerate}

The minimal implementation uses an orthonormal spinor Wannier basis, a fixed reference chemical potential, and a fixed local exchange-field vertex under the displacement. Thus the electronic mixed derivative contains the bubble term only. A direct mixed vertex is introduced as a separately controlled extension in \cref{app:direct}. Projector motion under displacement is outside the minimal approximation.

The magnon Hamiltonian is external. The required inputs are its positive energies, local spin frames, spin lengths, Nambu ordering, and paraunitary eigenvectors. The present formalism does not recompute exchange tensors or magnon bands. It also does not assert an operator identity between a KKR inverse-scattering-matrix displacement and a DFPT Hamiltonian derivative; only the common Dyson-insertion structure is used.

\section{Spinor Wannier Hamiltonian and exchange field}
\label{sec:spinor-model}

\subsection{Wannier representation}

Let $|w_{n\rr}\rangle$ denote orthonormal spinor Wannier functions, and adopt atomic-position Bloch sums
\begin{equation}
|w_{n\kk}\rangle
=\frac{1}{\sqrt N}\sum_{\rr}
 e^{\ii\kk\cdot(\rr+\bm\tau_n)}|w_{n\rr}\rangle.
\label{eq:atomic-bloch}
\end{equation}
Writing
\begin{equation}
H_{nm}(\rr)=\langle w_{n\mathbf 0}|H|w_{m\rr}\rangle,
\end{equation}
the Hamiltonian matrix in this gauge is
\begin{equation}
H_{nm}(\kk)=\sum_{\rr}
 e^{\ii\kk\cdot(\rr+\bm\tau_m-\bm\tau_n)}H_{nm}(\rr).
\label{eq:atomic-hamiltonian}
\end{equation}
The retarded Green function is
\begin{equation}
G^R_{\kk}(\varepsilon)=[(\varepsilon+\ii\eta)I-H(\kk)]^{-1}.
\label{eq:green}
\end{equation}
All spin-orbit terms remain in $H$ and hence in $G^R$. Equation~\eqref{eq:atomic-hamiltonian} also fixes the reciprocal-lattice sewing used below; no cell-gauge Fourier transform is mixed into the atomic-position convention.

\subsection{Time-reversal decomposition}

The relativistic localized-orbital formalism separates the time-reversal-symmetric Hamiltonian, which contains the scalar crystal Hamiltonian and SOC, from the time-reversal-broken exchange field~\cite{MartinezCarracedo2023}. We assume that no additional time-reversal-breaking one-particle term, such as an external orbital or Zeeman field, is present; otherwise that term must be separated explicitly. In a general Wannier gauge, let $B_\Theta(\kk)$ be the sewing matrix that represents time reversal between the bases at $-\kk$ and $\kk$. We define
\begin{equation}
H^\Theta(\kk)
=B_\Theta(\kk)H(-\kk)^*B_\Theta(\kk)^\dagger,
\label{eq:tr-transform}
\end{equation}
followed by
\begin{equation}
H_{\TRS}(\kk)=\frac{H(\kk)+H^\Theta(\kk)}{2},
\qquad
H_{\XC}(\kk)=\frac{H(\kk)-H^\Theta(\kk)}{2}.
\label{eq:tr-split}
\end{equation}
In a direct product of real orbital functions and spin eigenstates, $B_\Theta$ reduces to $I_{\rm orb}\otimes \ii\sigma_y$ up to the declared Wannier gauge. This reduction must not be assumed for a generic spinor Wannierization.

An explicit upstream decomposition $H=H_{\TRS}+H_{\XC}$ is preferable. For collinear spin blocks in a common orbital gauge, one may instead construct
\begin{equation}
H_{\rm avg}=\frac{H_\uparrow+H_\downarrow}{2},
\qquad
H_{\rm diff}=\frac{H_\uparrow-H_\downarrow}{2},
\label{eq:collinear-split}
\end{equation}
and lift $H_{\rm diff}$ along the chosen reference spin direction. Every route must reproduce $H$, yield a time-reversal-even $H_{\TRS}$ and a time-reversal-odd $H_{\XC}$, and state its factor convention.

\section{Local magnetic subspace and spin vertices}
\label{sec:local-torque}

\subsection{Hermitian local partition}

Let $P_\ell$ project onto the explicitly selected localized magnetic Wannier orbitals of site $\ell$, with $P_M=\sum_\ell P_\ell$. The orthonormal limit of the Hermitian local operator proposed in Ref.~\cite{MartinezCarracedo2023} is
\begin{equation}
\bL_\ell[X]=\frac{1}{2}(P_\ell X+XP_\ell).
\label{eq:local-partition}
\end{equation}
The local exchange fields are
\begin{equation}
H_{\XC,\ell}=\bL_\ell[H_{\XC}],
\qquad
\sum_\ell H_{\XC,\ell}
=\frac12(P_MH_{\XC}+H_{\XC}P_M).
\label{eq:local-xc}
\end{equation}
The difference between $H_{\XC}$ and the right-hand side is an explicit magnetic-support residual. The onsite-only choice $P_\ell H_{\XC}P_\ell$ is useful as a sensitivity test but is not the default.

The definition of $P_\ell$ is physical rather than merely numerical. Ref.~\cite{MartinezCarracedo2023} reports large changes when nonlocalized orbitals are included in the rotated subspace. A production calculation must therefore list the selected Wannier functions and compare at least a localized magnetic-orbital set with a larger atom-centered set.

\subsection{Local frames and transverse coordinates}

At each magnetic site, choose a right-handed frame
\begin{equation}
(\bm t_{\ell1},\bm t_{\ell2},\bm n_\ell),
\qquad
\bm t_{\ell1}\times\bm t_{\ell2}=\bm n_\ell,
\label{eq:local-frame}
\end{equation}
where $\bm n_\ell$ is the equilibrium moment direction. A small unit-vector fluctuation is
\begin{equation}
\bm e_\ell=\bm n_\ell
+\pi_{\ell1}\bm t_{\ell1}
+\pi_{\ell2}\bm t_{\ell2}+O(\pi^2).
\label{eq:pi-coordinate}
\end{equation}
The transverse-direction vertex is
\begin{equation}
\bT_{\ell a}
=\left.\frac{\partial H}{\partial\pi_{\ell a}}\right|_0.
\label{eq:direction-vertex}
\end{equation}
For a rigid local exchange field
\begin{equation}
H_{\XC,\ell}
=\frac12\Delta_\ell\otimes
(\bm n_\ell\cdot\bm\sigma),
\end{equation}
this becomes
\begin{equation}
\bT_{\ell a}
=\frac12\Delta_\ell\otimes
(\bm t_{\ell a}\cdot\bm\sigma).
\label{eq:rigid-direction-vertex}
\end{equation}

\subsection{Rotation-angle vertex and exchange-field-only rotation}

For an infinitesimal rotation around a unit axis $\bm u_{\ell c}$, the first-order exchange-field perturbation is~\cite{MartinezCarracedo2023}
\begin{equation}
\Gamma_{\ell c}
=\frac{\ii}{2}
[H_{\XC,\ell},\bm u_{\ell c}\cdot\bm\sigma]
=-\ii[S_{\ell c},H_{\XC,\ell}].
\label{eq:rotation-vertex}
\end{equation}
Only $H_{\XC,\ell}$ is rotated. Commuting the generator with the full Hamiltonian would rotate SOC relative to the lattice and does not implement the intended DFT-to-spin mapping~\cite{MartinezCarracedo2023}.

Direction and angle derivatives differ by a tangent-plane rotation. With $\bm u_{\ell1}=\bm t_{\ell1}$ and $\bm u_{\ell2}=\bm t_{\ell2}$,
\begin{equation}
\Gamma_{\ell1}=-\bT_{\ell2},
\qquad
\Gamma_{\ell2}=+\bT_{\ell1}.
\label{eq:coordinate-map}
\end{equation}
This distinction is essential for circular magnon combinations and torque signs.

\subsection{Finite-q representation of the local partition}

The real-space local partition generally produces a vertex that depends on both $\kk$ and $\qq$. Introduce the q-dependent site selector
\begin{equation}
P_\ell(\qq)
=\left\langle w_{\kk+\qq}\middle|
\sum_{\rr}e^{\ii\qq\cdot(\rr+\bm\tau_\ell)}P_{\rr\ell}
\middle|w_{\kk}\right\rangle.
\label{eq:q-projector}
\end{equation}
For diagonal orbital projectors in the atomic-position gauge,
\begin{equation}
[P_\ell(\qq)]_{nm}
=\delta_{nm}\delta_{s(n),\ell}
 e^{\ii\qq\cdot(\bm\tau_\ell-\bm\tau_n)}.
\label{eq:q-projector-atomic}
\end{equation}
Let $\Gamma_c^{\rm global}(\kk)$ be the uniform exchange-field rotation around the local axis selected for site $\ell$. Fourier transformation of the Hermitian local partition gives the project implementation rule
\begin{equation}
\Gamma_{\ell c}(\kk+\qq,\kk)
=\frac12\left[
P_\ell(\qq)\Gamma_c^{\rm global}(\kk)
+\Gamma_c^{\rm global}(\kk+\qq)P_\ell(\qq)
\right].
\label{eq:q-local-vertex}
\end{equation}
Only in the strictly onsite atomic-center limit is this matrix independent of $\kk$ and $\qq$. Equation~\eqref{eq:q-local-vertex} is the orthonormal finite-q implementation of the local-partition prescription and is classified as a project extension.

\section{Magnetic-force-theorem mixed derivative}
\label{sec:mixed-response}

At fixed chemical potential, the one-particle free-energy derivative can be written as
\begin{equation}
\frac{\partial\bF}{\partial\lambda}
=-\frac{1}{\pi}\operatorname{Im}
\int\dd\varepsilon\,f(\varepsilon)
\Tr[V_\lambda G^R],
\qquad
V_\lambda=\frac{\partial H}{\partial\lambda}.
\label{eq:first-derivative}
\end{equation}
Using
\begin{equation}
\frac{\partial G}{\partial\eta}=GV_\eta G,
\label{eq:dyson-derivative}
\end{equation}
we obtain
\begin{equation}
\frac{\partial^2\bF}{\partial\lambda\partial\eta}
=-\frac{1}{\pi}\operatorname{Im}
\int\dd\varepsilon\,f(\varepsilon)
\Tr\left[V_{\lambda\eta}G^R+V_\lambda G^R V_\eta G^R\right].
\label{eq:mixed-general}
\end{equation}
The second term is the mixed bubble. The same perturbative structure appears in the site-diagonal spin-lattice torque derivation of Ref.~\cite{Mankovsky2023SLC}; there the propagator and lattice vertex are KKR objects.

For a DFPT displacement,
\begin{equation}
g_{\kappa\mu}(\kk,\qq)
=\left\langle w_{\kk+\qq}\middle|
\frac{\partial H}{\partial u_{\kappa\mu}(\qq)}
\middle|w_{\kk}\right\rangle.
\label{eq:dfpt-g}
\end{equation}
This is the screened Hamiltonian derivative supplied by DFPT and transformed into the same Wannier gauge as $H$. The bubble-only approximation sets the mixed Hamiltonian derivative $V_{\lambda\eta}$ to zero. This is the minimal independent-perturbation implementation; the direct term is discussed in \cref{app:direct}.

The target energy coupling is
\begin{equation}
\mathcal H_{\rm SL}^{(1,1)}
=\sum_{\qq,\ell a,\kappa\mu}
\pi_{\ell a}(-\qq)
K_{\ell a,\kappa\mu}(\qq)
 u_{\kappa\mu}(\qq).
\label{eq:coupling-energy}
\end{equation}
The corresponding tangent effective field is $H^{\rm eff}_{\ell a}=-K_{\ell a,\kappa\mu}u_{\kappa\mu}$, while the physical torque is $\bm\tau_\ell=\bm n_\ell\times\bm H^{\rm eff}_\ell$, consistent with the torque-energy relation summarized in Ref.~\cite{Mankovsky2023SLC}.

\section{Finite-q Wannier kernel and gauge covariance}
\label{sec:finite-q}

Using the Fourier conventions in \cref{eq:atomic-bloch}, the DFPT vertex transfers an electron from $\kk$ to $\kk+\qq$, while the spin vertex at $-\qq$ closes the loop. Define the complex retarded loop
\begin{equation}
\begin{split}
\bA^R_{\ell a,\kappa\mu}(\qq)
=&\frac{1}{N_k}\int\dd\varepsilon\,f(\varepsilon)
\sum_{\kk}\tr\Big[
\bT_{\ell a}(-\qq;\kk,\kk+\qq)
G^R_{\kk+\qq}(\varepsilon)\\
&\hspace{35mm}\times
g_{\kappa\mu}(\kk,\qq)
G^R_{\kk}(\varepsilon)\Big].
\end{split}
\label{eq:retarded-loop}
\end{equation}
This is the q-resolved Wannier realization proposed here. The phonon-like effective field of Ref.~\cite{Mankovsky2023SLC} motivates the finite-q object, but the explicit momentum-space DFPT loop is a project extension.

A real-space Hermitian response need not have a real Fourier coefficient at a generic q. The physical kernel is obtained from the retarded/advanced q pair,
\begin{equation}
K^{\rm bub}_{\ell a,\kappa\mu}(\qq)
=-\frac{1}{2\pi\ii}
\left[
\bA^R_{\ell a,\kappa\mu}(\qq)
-\bA^R_{\ell a,\kappa\mu}(-\qq)^*
\right].
\label{eq:qpair-kernel}
\end{equation}
It follows that
\begin{equation}
K(-\qq)=K(\qq)^*.
\label{eq:q-reality}
\end{equation}
At a self-inverse q point, \cref{eq:qpair-kernel} reduces to the familiar $-\operatorname{Im}\bA^R/\pi$. Applying an elementwise imaginary part to one generic-q loop would incorrectly discard the phase of the Fourier coefficient.

If $\kk+\qq=\overline{\kk+\qq}+\bm G$, the atomic-position gauge uses
\begin{equation}
D_{\bm G}=\operatorname{diag}_n
 e^{\ii\bm G\cdot\bm\tau_n}\otimes I_{\rm spin},
\end{equation}
and
\begin{equation}
G(\kk+\qq)=D_{\bm G}^\dagger
G(\overline{\kk+\qq})D_{\bm G}.
\label{eq:wrap-green}
\end{equation}
The final-state index of $g(\kk,\qq)$ transforms consistently. After conversion, Hermitian lattice perturbations obey
\begin{equation}
g(\kk,\qq)^\dagger=g(\kk+\qq,-\qq),
\label{eq:g-hermiticity}
\end{equation}
with the appropriate wrapped indices.

\section{Projection onto phonon modes}
\label{sec:phonons}

For a phonon branch $\nu$, let $e_{\kappa\mu,\nu}(\qq)$ denote the eigenvector in a declared normalization. The zero-point displacement is
\begin{equation}
\xi_{\kappa\mu,\nu}(\qq)
=\sqrt{\frac{\hbar}{2M_\kappa\omega_{\qq\nu}}}
 e_{\kappa\mu,\nu}(\qq).
\label{eq:zpm}
\end{equation}
The Cartesian torque kernel projects to
\begin{equation}
V_{\ell a,\nu}(\qq)
=\sum_{\kappa\mu}
K_{\ell a,\kappa\mu}(\qq)
\xi_{\kappa\mu,\nu}(\qq).
\label{eq:phonon-proj}
\end{equation}
If the electron-phonon vertex is already mode normalized,
\begin{equation}
g_\nu(\kk,\qq)
=\sum_{\kappa\mu}\xi_{\kappa\mu,\nu}(\qq)
 g_{\kappa\mu}(\kk,\qq),
\label{eq:mode-g}
\end{equation}
then \cref{eq:zpm} must not be applied again.

Near an acoustic mode at $\qq=0$, the factor $1/\sqrt{\omega}$ makes a translational-sum-rule error visible. The Cartesian response to a rigid translation must therefore be checked before phonon projection. An exactly zero or imaginary frequency is reported rather than regularized silently.

\section{Projection onto external magnon modes}
\label{sec:magnons}

Let $N_m$ be the number of magnetic sublattices. Define the magnon Nambu spinor
\begin{equation}
\Psi_m(\qq)=
\begin{pmatrix}\bm a_{\qq}\\ \bm a_{-\qq}^\dagger\end{pmatrix}
=T_m(\qq)\Gamma_m(\qq),
\label{eq:magnon-nambu}
\end{equation}
where the external transformation satisfies
\begin{equation}
T_m^\dagger(\qq)\Sigma_m T_m(\qq)=\Sigma_m,
\qquad
\Sigma_m=\operatorname{diag}(I_{N_m},-I_{N_m}).
\label{eq:paraunitary}
\end{equation}
This is the standard paraunitary structure of a stable quadratic boson Hamiltonian~\cite{Colpa1978}.

In the local spin frames, the linearized unit-spin coordinates are
\begin{align}
\pi_{\ell1}(-\qq)
&=\frac{a_{\ell,-\qq}+a_{\ell,\qq}^\dagger}{\sqrt{2S_\ell}},\\
\pi_{\ell2}(-\qq)
&=\frac{-\ii a_{\ell,-\qq}+\ii a_{\ell,\qq}^\dagger}{\sqrt{2S_\ell}}.
\label{eq:hp-linear}
\end{align}
Collect these coefficients into $C_m$,
\begin{equation}
\bm\pi^T(-\qq)=\Psi_m^\dagger(\qq)C_m.
\label{eq:cm-map}
\end{equation}
For a phonon Nambu spinor $\Psi_p$ and coordinate matrix $C_p$, the coupling in the original boson bases is
\begin{equation}
V_{mp}(\qq)=C_m V_{\pi\nu}(\qq)C_p.
\label{eq:nambu-coupling}
\end{equation}
Substitution of $\Psi_m=T_m\Gamma_m$ and $\Psi_p=T_p\Gamma_p$ gives
\begin{equation}
\widetilde V_{mp}(\qq)
=T_m(\qq)^\dagger V_{mp}(\qq)T_p(\qq).
\label{eq:mode-transform}
\end{equation}
No additional metric occurs in \cref{eq:mode-transform}; the metric belongs to the eigenproblem and normalization in \cref{eq:paraunitary}. Adding it again would change the operator coefficients incorrectly.

The positive-frequency rows and particle/hole phonon columns of $\widetilde V_{mp}$ define normal and anomalous couplings. For a one-sublattice ferromagnet with no anomalous magnon mixing,
\begin{equation}
g^{mp}_\nu(\qq)
=\frac{V_{1\nu}(\qq)+\ii V_{2\nu}(\qq)}{\sqrt{2S}}.
\label{eq:fm-circular}
\end{equation}
In a rotating-wave treatment, an isolated set of magnon and phonon bands is coupled by
\begin{equation}
\mathcal H_{\rm RWA}(\qq)=
\begin{pmatrix}
\Omega_m(\qq) & g^{mp}(\qq)\\
[g^{mp}(\qq)]^\dagger & \Omega_{ph}(\qq)
\end{pmatrix}.
\label{eq:rwa}
\end{equation}
A full bosonic BdG assembly retains the anomalous blocks. Examples of topological magnon-polaron BdG models can be found in Ref.~\cite{Klogetvedt2023}; the present work only supplies the microscopic coupling input.

\section{Validation hierarchy}
\label{sec:validation}

The method is sensitive to basis localization, spin conventions, and q-space phases. Validation is therefore part of the definition of the calculation.

\subsection{Exchange decomposition and locality}

The time-reversal decomposition must reconstruct $H$, while $H_{\TRS}$ and $H_{\XC}$ have the correct parity. The local fields and torque vertices must be Hermitian. The local partition must close on the magnetic support in \cref{eq:local-xc}. The norms of the residual support and the difference between local and onsite projections are reported. Following the sensitivity found in Ref.~\cite{MartinezCarracedo2023}, the result is compared for a localized magnetic-orbital set and a larger atom-centered set.

\subsection{Spin-vertex checks}

A central finite rotation of $H_{\XC,\ell}$ must reproduce \cref{eq:rotation-vertex}. A rotation about the local moment direction must give zero. Scaling SOC changes the Green function but not a vertex generated from a fixed $H_{\XC}$. Direction and angle vertices must satisfy \cref{eq:coordinate-map}.

\subsection{Mixed finite difference}

Within a linearized Wannier model, the analytic mixed response is compared with
\begin{equation}
K^{\rm FD}=
\frac{\bF(+\delta\pi,+\delta u)-\bF(+\delta\pi,-\delta u)
-\bF(-\delta\pi,+\delta u)+\bF(-\delta\pi,-\delta u)}
{4\delta\pi\delta u}.
\label{eq:mixed-fd}
\end{equation}
The same fixed-vertex approximation must be used on both sides. For a collinear stationary reference with spin-rotation-invariant electronic terms, the linear transverse kernel should vanish when SOC and other anisotropic spin interactions are removed.

\subsection{Gauge, translation, and bosonic invariance}

The DFPT Hermiticity relation \cref{eq:g-hermiticity}, the q reality condition \cref{eq:q-reality}, and atomic-gauge wrapping are tested directly. At $\qq=0$, the sum of Cartesian kernels over a rigid translation of all atoms must vanish before phonon normalization. Independent phase changes of phonon and magnon eigenvectors may change individual complex couplings but must leave $|g|$, degenerate-subspace singular values, and polaron eigenvalues invariant.

\section{Implementation contract}
\label{sec:implementation}

The numerical implementation is organized into four layers.

\begin{enumerate}
\item \textbf{Electronic model.} Load and canonicalize the spinor Wannier Hamiltonian, time-reversal metadata, exchange-field decomposition, magnetic projectors, and local frames.
\item \textbf{Mixed electronic response.} Stream screened DFPT perturbations, build wrapped Green functions, accumulate the complex retarded loop, and complete q pairs.
\item \textbf{Mode projection.} Convert Cartesian kernels to phonon zero-point modes and then to external magnon quasiparticles.
\item \textbf{Validation and output.} Store HDF5 metadata, source hashes, restart status, decomposition residuals, q-symmetry tests, subspace sensitivity, and convergence tables.
\end{enumerate}

The reference matrix contraction for a batch of torque vertices $T_t$ and displacement perturbations $g_p$ uses
\begin{equation}
A_p=G_{\kk+\qq}g_pG_\kk,
\qquad
L_{tp}=\tr[T_tA_p],
\label{eq:batch-contraction}
\end{equation}
without materializing a tensor carrying both vertex labels and two matrix indices. q pairs are the primary parallel unit, while deterministic serial reductions remain available for validation.

All input files declare Bloch and Fourier gauges, basis ordering, time-reversal sewing convention, magnetic orbital masks, site projection, spin coordinate, and DFPT normalization. Missing metadata is an error rather than an invitation to infer conventions from array shapes.

\section{Discussion and limitations}
\label{sec:discussion}

The construction preserves the main physical distinction required by relativistic LKAG: the spin degree of freedom is rotated relative to a fixed crystal. SOC therefore remains in the propagator, while the local perturbation is generated from the exchange field. This avoids interpreting a global spinor rotation of the complete Hamiltonian as a physical moment rotation.

The local torque kernel should not automatically be called a pure single-ion anisotropy derivative. As emphasized by the spin-lattice source~\cite{Mankovsky2023SLC}, an effective field at one site can contain nonlocal anisotropic intersite contributions. The full Green functions similarly allow the Wannier kernel to collect propagation through the entire crystal. Direct projection onto magnon eigenvectors is therefore natural when the objective is the total linear hybridization.

Several approximations remain visible. First, a collinear lift of DFPT data may omit spin-flip lattice perturbations and $\partial H_{\rm SOC}/\partial u$. Second, the bubble-only approximation omits the displacement dependence of the exchange-field torque vertex. Third, fixed reference projectors omit Wannier-gauge or Pulay-like projector motion. Fourth, the magnetic-force-theorem derivative is evaluated at fixed chemical potential in the minimal implementation. These approximations can be relaxed independently if the required upstream data are available.

The methodological novelty is correspondingly narrow and testable: a q-resolved spinor-Wannier implementation of displacement-induced magnetic torque using screened DFPT matrix elements, completed by gauge-controlled projection onto independently calculated phonon and magnon modes. The local relativistic torque formalism, the KKR spin-lattice expression, DFPT, Wannier interpolation, and paraunitary boson theory are established ingredients and should be cited as such.

\section{Conclusion}
\label{sec:conclusion}

We have specified a relativistic mixed spin-lattice response suitable for a Wannier-based implementation and direct magnon-polaron construction. The exchange field is separated from the SOC-containing time-reversal-symmetric Hamiltonian, partitioned over explicit localized magnetic subspaces, and differentiated with respect to local transverse spin coordinates. A screened DFPT displacement vertex enters the Green-function bubble through $\partial_uG=GgG$. General finite-q coefficients are retained as complex quantities and obtained from retarded q pairs. The Cartesian response is projected first onto phonon zero-point displacements and then onto external paraunitary magnon eigenvectors without an additional bosonic metric.

The resulting implementation has a clear approximation hierarchy: bubble-only response as the first validated milestone, followed by an optional exchange-field-resolved direct mixed vertex. The required finite-difference, time-reversal, local-projection, translation, gauge, subspace, and bosonic rephasing tests make the formalism operational rather than merely symbolic.

\appendix
\section{Rotation angles, direction coordinates, and physical torque}
\label{app:coordinates}

For a small rotation vector $\delta\bm\theta_\ell$,
\begin{equation}
\delta\bm e_\ell=\delta\bm\theta_\ell\times\bm n_\ell.
\end{equation}
Choosing $\delta\bm\theta=\delta\theta_1\bm t_1+\delta\theta_2\bm t_2$ gives
\begin{equation}
\delta\bm e=-\delta\theta_1\bm t_2+\delta\theta_2\bm t_1,
\end{equation}
which proves \cref{eq:coordinate-map}. If $\bm H^{\rm eff}=-\partial\bF/\partial\bm e$, then
\begin{equation}
\delta\bF=-\delta\bm\theta\cdot
(\bm n\times\bm H^{\rm eff}),
\end{equation}
so the physical torque is $\bm\tau=\bm n\times\bm H^{\rm eff}$ and a rotation-angle component satisfies $\tau_c=-\partial\bF/\partial\theta_c$. The direction-coordinate kernel must therefore be rotated in the tangent plane before it is interpreted as a torque-vector response.

\section{Retarded/advanced completion at finite q}
\label{app:qpair}

Let $V_A(-\qq)$ and $V_B(\qq)$ be Hermitian-conjugate Fourier components of two real-space perturbations. The retarded loop is
\begin{equation}
A^R_{AB}(\qq)=\int\tr[V_A(-\qq)G^R_{\kk+\qq}V_B(\qq)G^R_\kk].
\end{equation}
Taking the complex conjugate of the retarded loop at $-\qq$, using $V_X(\qq)^\dagger=V_X(-\qq)$, cyclically permuting the trace, and shifting the dummy momentum by $\qq$ gives the advanced loop with the same ordered pair of perturbations,
\begin{equation}
A^A_{AB}(\qq)=A^R_{AB}(-\qq)^*.
\label{eq:qpair-advanced}
\end{equation}
This step is important for two different vertices: the momentum shift, rather than an assumed elementwise reality of $A^R_{AB}(\qq)$, restores the original $A,B$ ordering.
The spectral discontinuity is therefore
\begin{equation}
-\frac{A^R_{AB}(\qq)-A^A_{AB}(\qq)}{2\pi\ii},
\end{equation}
which yields \cref{eq:qpair-kernel}. At a self-inverse q point, $A^A=A^{R*}$ and the expression reduces to $-\operatorname{Im}A^R/\pi$.

\section{Direct mixed vertex and moving projectors}
\label{app:direct}

The exact one-particle mixed derivative contains
\begin{equation}
\bT^{(1)}_{\ell a,\kappa\mu}
=\frac{\partial^2H}
{\partial\pi_{\ell a}\partial u_{\kappa\mu}},
\label{eq:direct-vertex}
\end{equation}
in addition to the bubble. If the displacement derivative of the exchange field, $g_{\XC}=\partial H_{\XC}/\partial u$, is available and the projectors are fixed, then a rotation-coordinate version is
\begin{equation}
\Gamma^{(1)}_{\ell c,\kappa\mu}
=\frac{\ii}{2}
[\bL_\ell(g_{\XC,\kappa\mu}),
 \bm u_{\ell c}\cdot\bm\sigma].
\label{eq:direct-rotation}
\end{equation}
The full DFPT vertex must not replace $g_{\XC}$ in this commutator, because it may contain scalar and SOC derivatives that are fixed to the lattice.

If the magnetic projector itself changes under displacement,
\begin{equation}
\partial_u\bL_\ell[X]
=\frac12[(\partial_uP_\ell)X+X(\partial_uP_\ell)]
+\bL_\ell[\partial_uX].
\label{eq:projector-motion}
\end{equation}
The first term is a projector-response correction and requires additional gauge information. Consequently, the direct term is enabled only when the exchange-field derivative and projector policy are supplied explicitly and validated by a mixed finite difference.


\bibliographystyle{apsrev4-2}
\bibliography{references}
\end{document}
